% S02 continuation: printed pages 44–48, physical PDF133–137.
% Historical source: Nallino, Pars I.
% TRANSCRIPTION CANDIDATE — compilation and final visual QA not certified.
%
% Body responsibility: NALLINO_TRANSLATION
% Historical notes: NALLINO_APPARATUS
% Editorial checking register: MODERN_EDITORIAL
%
% Controlling source SHA-256:
% 544c16b6355c9b74e281aded657d657224bff738366d5260e1a610d31b0d6297
%
% Requires the original master PDF in the same directory.
% Compile with LuaLaTeX.
%
% Source paragraph, note, and figure anchors are supplied.
% No new target-language translation is included.

\documentclass[10pt]{article}
\usepackage[
  paperwidth=220mm,
  paperheight=320mm,
  left=17mm,
  right=17mm,
  top=19mm,
  bottom=19mm
]{geometry}

\usepackage{fontspec}
\usepackage{polyglossia}
\setdefaultlanguage{latin}
\setotherlanguages{arabic,greek}

\IfFontExistsTF{Libertinus Serif}
  {\setmainfont{Libertinus Serif}}
  {\IfFontExistsTF{FreeSerif}
    {\setmainfont{FreeSerif}}
    {\setmainfont{DejaVu Serif}}}

\IfFontExistsTF{Amiri}
  {\newfontfamily\arabicfont[Script=Arabic]{Amiri}}
  {\newfontfamily\arabicfont[Script=Arabic]{DejaVu Sans}}

\IfFontExistsTF{FreeSerif}
  {\newfontfamily\greekfont{FreeSerif}}
  {\newfontfamily\greekfont{DejaVu Serif}}

\usepackage{amsmath}
\usepackage{graphicx}
\usepackage{multicol}
\usepackage{hyperref}
\hypersetup{hidelinks}

\setlength{\parindent}{1em}
\setlength{\parskip}{3pt}
\setlength{\columnsep}{6mm}
\emergencystretch=1em

\newcommand{\MasterPDF}
  {30_NALLINO_PARS_I_II_III_MASTER_1162P.pdf}

\newcommand{\A}[1]{\hypertarget{#1}{}}
\newcommand{\NM}[1]{\textsuperscript{(#1)}}
\newcommand{\NN}[2]{\par\noindent\textup{(#1)}\ #2}
\newcommand{\AR}[1]{\textarabic{#1}}
\newcommand{\GR}[1]{\textgreek{#1}}
\newcommand{\U}[1]{\textsuperscript{\texttt{[#1]}}}

\newcommand{\SP}[3]{%
  \clearpage
  \A{AB01-PDF#1}%
  \begin{center}\textbf{#2\quad #3}\end{center}%
}

% Crop a source region without redrawing it.
%
% Arguments:
% #1 = physical PDF page, one-indexed
% #2 = x0, points from left
% #3 = y0, points from top
% #4 = x1, points from left
% #5 = y1, points from top
% #6 = desired display width
%
% Native PDF page dimensions are measured rather than assumed.
% The source PDF itself remains the untouched raw witness.

\newsavebox{\SourcePageBox}
\newlength{\SourceBottomTrim}
\newlength{\SourceRightTrim}

\newcommand{\SourceCrop}[6]{%
  \begingroup
  \sbox{\SourcePageBox}{%
    \includegraphics[page=#1]{\MasterPDF}%
  }%
  \setlength{\SourceBottomTrim}{%
    \dimexpr\ht\SourcePageBox+\dp\SourcePageBox-#5bp\relax
  }%
  \setlength{\SourceRightTrim}{%
    \dimexpr\wd\SourcePageBox-#4bp\relax
  }%
  \includegraphics[
    page=#1,
    trim={#2bp \SourceBottomTrim \SourceRightTrim #3bp},
    clip,
    width=#6
  ]{\MasterPDF}%
  \endgroup
}

\begin{document}

% ================================================================
% PHYSICAL PDF133 — PRINTED PAGE 44
% ================================================================

\SP{0133}{44}{C. A. NALLINO}

\A{AB01-PDF0133-P01}
partem quam Sol medio suo itinere ab initio Arietis ad initium
Librae percurrit, scilicet, ut supra vidimus,
$183^\circ56'12''$. Arcus KLRM, semicircumferentia excentrici,
$180^\circ$ continet; quisque igitur arcuum PK, YM est dimidia
pars eorum $3^\circ56'12''$, quibus Sol medio suo itinere
$180^\circ$ superat, idest $1^\circ58'6''$. — Manifestum etiam
est arcum PKLR\NM{1} excentrici id esse quod Sol, medio suo
itinere, ab initio Arietis ad initium Cancri percurrit; igitur
$92^\circ14'10''$ erit; et quia arcus PKL, ob ea quae nuper
diximus, $91^\circ58'6''$ complectitur, arcus LR $0^\circ16'4''$
continebit. Item manifestum est perpendicularem PQ sinum esse
arcus PK, et perpendicularem RG sinum arcus LR; igitur PQ fere
$2^{p}3'39''$, et RG fere $0^{p}16'45''$ erit. Quoniam linea KM
parallela est lineae AC, erit linea ES aequalis PQ; et similiter,
quia LN parallela est lineae BD, erit FS aequalis RG. Est itaque
etiam latus EF trianguli rectanguli ESF notum; quadratum lateris
ES est fere $4^{p}14'48''$, quadratum lateris FS fere
$0^{p}4'41''$, summa igitur $4^{p}19'29''$\NM{2} efficit quadratum
lateris EF, cuius radix $2^{p}4'\,3/4$ est quantitas lineae EF
ambo centra coniungentis. Ea ratione autem, qua quadrans circuli
triangulum rectangulum ESF circumdantis $90^\circ$ complectitur,
et semidiametrus $60^{p}$, erit arcus EF circiter
$1^\circ59'$\NM{3}; et haec est maxima anomalia motus Solis per
has observationes manifesta \textit{(b)}.

\A{AB01-PDF0133-P02}
Nunc in quantitatem arcus BH eclipticae inquiramus; qua cognita,
etiam arcus HA reliquus notus erit. Punctum T est locus apogei in
circulo Solis excentrico; nam, si lineam EF, ambo centra
coniungentem, usque ad eclipticam producimus, ipsa circulum KLMN
in T, et circulum eclipticae in H secat. Necesse est ut rationem
lineae EF ad semidiametrum EH cognoscamus, nec non quantitatem
arcus BH eclipticae. Iam vidimus lineam EF $2^{p}4'\,3/4$
complecti ea ratione qua semidiametrus $60^{p}$ continet; eadem
ratione igitur linea EH, quae, ut EB, semidiametrus est, lineam
EF 28 vicibus et $5/6$ fere continebit. Quoniam linea FS, ut
diximus, [$0^{p}16'45''$] est, si linea EF $60^{p}$ ponatur,
erit iuxta hanc rationem linea FS circiter $8^{p}4'$;

\par\noindent
\begin{minipage}[t]{.43\textwidth}
  \vspace{0pt}
  \centering
  \A{AB01-PDF0133-F01}
  % Raw source region, top-origin PDF coordinates:
  % x0=108, y0=436, x1=313, y1=633.
  \SourceCrop{133}{108}{436}{313}{633}{\linewidth}
\end{minipage}
\hfill
\begin{minipage}[t]{.54\textwidth}
  \vspace{0pt}
  hoc enim invenitur si $28\,5/6$ vicibus [$0^{p}16'45''$]
  multiplicantur. — Aut, si mavis, lineam FS in semidiametrum
  EH multiplica; invenies $16^{p}45'$, quae, per lineam EF
  scilicet per $2^{p}4'\,3/4$ divisa, dabunt $8^{p}4'$, quod
  est [sinus] anguli BEH; arcus igitur BH $7^\circ43'$ circiter
  continet. Itaque manifestum est punctum T apogei in
  excentrico, a puncto solstitii aestivi versus partem
  anteriorem signorum $7^\circ43'$ distare, scilicet
  $82^\circ17'$ ab initio Arietis esse. Hoc invenire volebamus.
  Observatio iuxta quam hunc calculum fecimus, anno 1194 aerae
  Dhū ’l-qarnayn\NM{4} fuit, quo iter Solis ab initio Arietis
  ad initium Cancri et ad initium Librae observavimus
  \textit{(c)}.
\end{minipage}

\A{AB01-PDF0133-P03}
Restat nunc nobis ut hanc anomaliam ad gradus eclipticae
distribuamus, et portiones eius quae ad singulos gradus
pertinent in tabulis

\begin{multicols}{2}
\footnotesize

\A{AB01-PDF0133-N01}
\NN{1}{Male codex PKL.}

\A{AB01-PDF0133-N02}
\NN{2}{Male secundae apud Platonem $59''$.}

\A{AB01-PDF0133-N03}
\NN{3}{Plato male $1^\circ58'$; Ibn Yūnus
(apud Caussin, \textit{Le livre de la grande table Hakémite},
in Notices et extraits des mss. de la Bibl. Nation., t. VII,
Paris 1804, p. 155) scribit: « al-Battāni aequationem totam
Solis calculavit, invenitque $1^\circ59'10''$ esse ».
Halley p. 917, et Delambre, \textit{Astr. moyen âge}, p. 36,
qui hodiernis tabulis trigonometricis usi sunt, invenere
quoque $1^\circ59'10''$. Codicis lectio igitur bona est; in
tabulis aequationis Solis (fol. 191,r.) re vera maxima
aequatio est $1^\circ59'10''$.}

\A{AB01-PDF0133-N04}
\NN{4}{Incipit, iuxta usum auctoris nostri, die 1 Septembris
882 post Chr. n.}

\end{multicols}

% ================================================================
% PHYSICAL PDF134 — PRINTED PAGE 45
% ================================================================

\SP{0134}{45}{AL-BATTANI OPUS ASTRONOMICUM}

\A{AB01-PDF0134-P01}
describamus, ut facile possit aequatio motus Solis inveniri.
Ostendit Ptolemaeus\NM{1} motus varios duobus modis effingi
posse; quorum alter est ut ponatur planetam habere circulum
eclipticae concentricum, super quem alius circulus sit, cuius
centrum circumferentiam primi percurrat. Hic secundus circulus
est circulus minor, Terram non ambiens, cuius centrum circulus
maior rotare facit versus partem successionis signorum iuxta
quantitatem motus longitudinalis planetae versus partem in
quam signa succedunt. Sidus ipsum in epicyclo, qui est circulus
minor, versus praecedentem aut versus subsequentem partem
movetur; aut circulus minor planetam versus alterutram partem
rotare facit. Hic est motus anomaliae quae ad planetam pertinet.
— Alter explicandi modus est ut ponatur sidus habere circulum
concentricum eclipticae, et alterum excentricum sed illi
aequalem, qui eum duobus locis secet Planeta est in excentrico,
sive hic eum movet, sive ipse in excentrico movetur.
— Utraque ratio anomaliam pariter explicat; nos a prima
incipiemus.

% No punctuation has been supplied between "secet" and "Planeta".

\A{AB01-PDF0134-P02}
Circulum ABCD eclipticae circa centrum E describamus; sit A,
initio, centrum epicycli HF\NM{2}. Diametrum AC usque ad H
producamus, quod est punctum apogei in epicyclo; sit F locus
Solis in epicyclo, et ab eo perpendicularem supra lineam AH
usque ad punctum M ducamus; denique lineam AF signemus,
lineae AH aequalem, nam ambae sunt semidiametri epicycli.
Ob ea quae in hoc ipso capite iam diximus, patet eam
$2^{p}4'\,3/4$ complecti, ut semidiametrum EF epicycli in
figura praecedenti. Quo cognito, motum Solis in epicyclo versus
partem contrariam successioni signorum observa, vel quantitatem
qua epicyclus motu medio diurno Solem versus illam partem
movet, iuxta rationem qua circumferentia epicycli in
$360^\circ$ dividitur. Motus medius Solis, qui observationibus
invenitur, est motus centri epicycli versus partem successionis
signorum; et hic motus quoque ea ratione ponitur qua circulus
ABCD $360^\circ$ complectitur.

\A{AB01-PDF0134-P03}
Arcus HF, qui est arcus epicycli inter locum Solis et punctum
apogei, $30^\circ$ contineat e gradibus quorum epicyclus 360
continet. Lineam EF in hac figura ducamus; et in arcum lineae
FM, qui est variatio motus Solis eo loco, inquiramus.
Linea EA, semidiametrus circuli concentrici eclipticae, 60
partes habet ex iis quarum diametrus AC 120 continet; linea
igitur EH, quae centrum circuli parecliptici\NM{3} cum apogeo
epicycli coniungit, $62^{p}4'45''$ erit. Et quia triangulum FMA
est rectangulum, quadratum lineae AF erit ut summa quadratorum
AM, FM; angulus autem MAF\NM{4} notus est, ergo etiam linea
FM est nota, a qua deinde co-

\begin{multicols}{2}
\footnotesize

\A{AB01-PDF0134-N01}
\NN{1}{\textit{Almag.} III, 3, 4 et 5
(ed. Halma, t. I, p. 170–183, 183–190, 190–199);
quoad Solem tamen hypothesim excentrici commodiorem censet.}

\A{AB01-PDF0134-N02}
\NN{2}{Cod. et Plato HFT; in eorum figura enim punctum T
ponitur loco intersectionis epicycli et circuli magni,
inter A et D. Sed ex iis quae sequuntur patet T esse punctum
circuli magni indicans directionem apogei.}

\A{AB01-PDF0134-N03}
\NN{3}{Scilicet paralleli seu concentrici eclipticae.
Cum nullum habeat nomen apud scriptores nostrates, ita verto
Arabicum \textit{al-mumaththal bi falak al-burūǵ}\U{C01}
« similem factum eclipticae » seu \textit{al-mumaththal}
tantum, qui apud Ptolemaeum interdum
\GR{ὁ ὁμόκεντρος τῷ κόσμῳ κύκλος}
« circulus mundo concentricus » vel
\GR{ὁ ὁμόκεντρος τῷ διὰ μέσων τῶν ζωδίων κύκλος}\U{C02}
« circulus concentricus circulo per medium signorum
transeunti » (i. e. eclipticae) vocatur. In astronomia
Syriaca Barhebraei vertenda, Nau eum
\textit{intersphère de la similitude} appellat; Golius,
in versione al-Farghāni, et Sédillot, in versione Ulugh Beg,
\textit{sphaeram} vel \textit{circulum homocentricum}.
— Est circulus concentricus eclipticae in cuius plano iacet;
descriptus est in superficie sphaerae Solis vel planetarum,
\textit{al-mumaththal} quoque vocata, quae est corpus solidum,
concentricum eclipticae, duabus superficiebus parallelis
contentum, quarum exterior superficiem exteriorem sphaerae
circuli excentrici (deferentis) in puncto apogeo excentrici
tangit; interior superficiem interiorem sphaerae excentrici
in perigeo. — Movetur ut signa zodiaci ob praecessionem
moventur, ab ortu ad occasum.}

\A{AB01-PDF0134-N04}
\NN{4}{Cod. HEF.}

\end{multicols}

% ================================================================
% PHYSICAL PDF135 — PRINTED PAGE 46
% ================================================================

\SP{0135}{46}{C. A. NALLINO}

\A{AB01-PDF0135-P01}
gnoscitur reliquum latus AM trianguli, quod est [cosinus]
arcus FH et anguli FAH\NM{1}. Igitur est etiam linea EM nota.
Cum triangulum FME rectangulum sit, eius hypotenusa EF
cognoscitur, ex qua linea FM deprehenditur; arcus autem super
FM est arcus anomaliae.

\A{AB01-PDF0135-P02}
Si, ut posuimus, arcus FH $30^\circ$ complectitur, eius sinus
erit 30 partium, e partibus quarum semidiametrus AF 60
continet; sed ea ratione qua linea AF $2^{p}4'45''$
complectitur, erit linea FM $1^{p}2'22''\,1/2$\NM{2}, reliqua
igitur AM $1^{p}48'2''$, et linea EM\NM{3}
$61^{p}48'2''$\NM{4}; unde patet [lineam EF] $61^{p}48'35''$
circiter esse \textit{(d)}. Sed ex partibus quarum linea EF
60 tantum habet, linea FM $1^{p}0'33''$\NM{5} continebit,
et arcus super eam positus $0^\circ57'49''$ circiter; et haec
est quantitas anomaliae motus Solis, scilicet arcus HF.
Est igitur arcus TA eclipticae $29^\circ2'11''$, qui antea
$30^\circ$ erat; nam epicycli centrum a T ad A motum est
ut Sol in epicyclo a H ad F.

\A{AB01-PDF0135-P03}
Nunc epicyclum GQP consideremus, cuius centrum est B; locum
Solis G ponamus, sitque arcus QG\NM{6}, quem Sol a puncto Q
apogei percurrit, $150^\circ$\NM{7}; remanent igitur
$30^\circ$ arcui PG, qui a loco Solis ad punctum perigei
porrigitur. Lineam EG et perpendicularem GK ducamus; patet
triangula BKG et GKE rectangula esse, et latera BG, BE haud
ignota; nam BG est semidiametrus epicycli\NM{8}, et BE
semidiametrus circuli eclipticae. Angulus et arcus GP dati
sunt; perpendicularis GK et reliqua linea BK notae;
[ergo etiam lineae KE et EG cognoscuntur]\NM{9}.

\par\noindent
\begin{minipage}[t]{.44\textwidth}
  \vspace{0pt}
  \centering
  \A{AB01-PDF0135-F01}
  % Source region: x0=120, y0=350, x1=329, y1=568.
  \SourceCrop{135}{120}{350}{329}{568}{\linewidth}
\end{minipage}
\hfill
\begin{minipage}[t]{.53\textwidth}
  \vspace{0pt}
  Cum arcus GP, ut posuimus, $30^\circ$ contineat, eius sinus
  quoque $30^{p}$ erit; arcus autem complementaris, qui super
  KB est, sinum habebit $51^{p}57'41''$\NM{10}.
  — At ex partibus quarum linea BG $2^{p}4'\,3/4$ habet,
  perpendicularis KG $1^{p}2'22''\,1/2$\NM{11} continebit;
  remanet linea KB $1^{p}48'2''$, qua re EK fere
  $58^{p}11'58''$, et EG fere $58^{p}12'34''$\NM{12}
  continebit \textit{(e)}. Sed ex partibus quarum EG 60
  habet, cathetus KG $1^{p}4'17''$\NM{13} complectetur,
  et arcus super eam positus $1^\circ1'24''$\NM{14},
  iuxta rationem qua circulus triangulum rectangulum BKG
  circumdans in $360^\circ$ dividitur. Et hic est arcus
  anomaliae, scilicet arcus PG. Qua re arcus NB\NM{15}
  eclipticae, quem cognoscere volebamus,
  $31^\circ1'24''$\NM{16} complectetur.
\end{minipage}

\begin{multicols}{2}
\footnotesize

\A{AB01-PDF0135-N01}
\NN{1}{Cod.: « quod est quod angulo FAH et arcui FH deest
ad quadrantem perficiendum ».}

\A{AB01-PDF0135-N02}
\NN{2}{Pro $22''$ Plato habet $55''$. Persaepe in Platonis
editione pro 2 legitur 5, qui error codicibus Latinis vel
editoris incuriae tribuendus est.}

\A{AB01-PDF0135-N03}
\NN{3}{Codex EF.}

\A{AB01-PDF0135-N04}
\NN{4}{Plato verba seqq. omittit et pro $2''$ habet $35''$.}

\A{AB01-PDF0135-N05}
\NN{5}{Plato $1^\circ33'$.}

\A{AB01-PDF0135-N06}
\NN{6}{Cod. QM.}

\A{AB01-PDF0135-N07}
\NN{7}{Male Plato 120.}

\A{AB01-PDF0135-N08}
\NN{8}{Cod.: eclipticae.}

\A{AB01-PDF0135-N09}
\NN{9}{Addidi, Platonem sequens; sed additio non est necessaria.}

\A{AB01-PDF0135-N10}
\NN{10}{Partes apud Plat. $59^{p}$.}

\A{AB01-PDF0135-N11}
\NN{11}{Deest $1/2$ in cod.; Plato pro $22''$ habet $55''$
(cfr. adnot. 2).}

\A{AB01-PDF0135-N12}
\NN{12}{Plato $58^{p}7'34''$.}

\A{AB01-PDF0135-N13}
\NN{13}{Secundae apud Platonem $13''$.}

\A{AB01-PDF0135-N14}
\NN{14}{Plato $1^\circ4'54''$.}

\A{AB01-PDF0135-N15}
\NN{15}{Cod. IB (\AR{ى} I pro \AR{ن} N)\U{C03}; Plato GB.
In figura codicis et Platonis punctum N lineaque NT desunt;
punctum T vero, ut supra notavi, perperam prope locum
intersectionis epicycli HF et eclipticae collocatum.}

\A{AB01-PDF0135-N16}
\NN{16}{Minuta in cod. $4'$; Plat. $1^\circ4'54''$.}

\end{multicols}

% ================================================================
% PHYSICAL PDF136 — PRINTED PAGE 47
% ================================================================

\SP{0136}{47}{AL-BATTANI OPUS ASTRONOMICUM}

\A{AB01-PDF0136-P01}
Dicit [auctor]: Nunc aliter haec explicabimus. Circa centrum E
et diametrum AC, sit circulus eclipticae ABC; circa centrum H,
excentricus FMG; et diametrus AC per ambo centra transeat.
Erit F punctum apogei in circulo parecliptico, G perigei.
Initio locum Solis in excentrico in M esse ponamus, et arcum FM
excentrici, quem Sol iam percurrit, $30^\circ$ complecti;
angulus FHM erit quoque $30^\circ$. Lineam EH, ambo centra
coniungentem, iam vidimus $2^{p}4'\,3/4$ esse.
Semidiametrum HM excentrici et lineam EM ducamus, HM usque
ad punctum L quo cum perpendiculari LE convenit protrahentes.
Triangulum HLE rectangulum est; eius angulus LHE aequat
angulum datum FHM; arcus super EL, ex circulo triangulum
HLE circumdanti, si circumferentia $360^\circ$ continet,
$30^\circ$ complectetur, eiusque sinus $30^{p}$ quoque ex
partibus quarum linea HE inter ambo centra 60 habet.
Remanebit linea LH, quadrantem perficiens, $51^{p}57'41''$,
nam arcus complementaris LH $60^\circ$ continet.
Sed ea ratione qua linea HE inter ambo centra $2^{p}4'\,3/4$
complectitur, linea EL $1^{p}2'22''\,1/2$\NM{1} erit, et linea
LH, quadrantem perficiens, $1^{p}48'2''$; itaque linea tota
LM $61^{p}48'2''$. Triangulum MLE rectangulum est;
eius hypotenusam EM $61^{p}48'35''$ esse vidimus.
Sed ratione qua lineae EM $60^{p}$ tribuuntur, EL est
$1^{p}33'$\NM{2}, et arcus super eam $0^\circ57'49''$,
cum circulus triangulum HLE circumdans $360^\circ$ sit.
Arcus igitur AB eclipticae $29^\circ2'11''$ fere remanebit.

\A{AB01-PDF0136-P02}
Item Solem in puncto D excentrici esse ponamus, arcumque
FD\NM{3} $150^\circ$ complecti; reliquus arcus DG, inter
locum Solis et perigeum, $30^\circ$ erit. Lineas EK et HD,
quarum utraque est sui circuli semidiametrus, ducamus;
item perpendicularem ES. Triangulum HSE rectangulum est;
latus EH, ambo centra coniungens, latus ES, angulus
DHG\NM{4} nota sunt; erunt ergo nota latus HS et angulus
HES\NM{5}. Cognoscuntur igitur quoque linea DG et hypotenusa
ED trianguli rectanguli ESD. Cum vero datus arcus DG et
datus angulus GHD $30^\circ$ sint ut diximus, eorumque sinus
$30^{p}$ pariter, arcus quoque ES circuli triangulum
rectangulum ESH circumdantis complectetur $30^\circ$ ex
partibus quarum totus ille circulus $360^\circ$ habet,
et eius sinus, nempe cathetus ES, $30^{p}$ quoque e
partibus quarum linea EH, semidiametrus huius circuli,
$60^{p}$ complectitur. Sed ratione qua lineae EH
$2^{p}4'\,3/4$ tribuuntur, cathetus ES erit
$1^{p}2'22''\,1/2$\NM{6}; qua re latus SH $1^{p}48'2''$
remanebit.

\par\noindent
\begin{minipage}[t]{.51\textwidth}
  \vspace{0pt}
  \centering
  \A{AB01-PDF0136-F01}
  % Source region: x0=65, y0=480, x1=298, y1=645.
  \SourceCrop{136}{65}{480}{298}{645}{\linewidth}
\end{minipage}
\hfill
\begin{minipage}[t]{.46\textwidth}
  \vspace{0pt}
  Semidiametrus HD excentrici $60^{p}$ continet, de quibus
  si SH dematur, remanent lineae SD $58^{p}11'58''$;
  igitur hypotenusa ED trianguli rectanguli ESD erit
  circiter $58^{p}12'32''$\NM{7}. Sed ex partibus quarum
  linea ED 60 habet, cathetus ES $1^{p}4'17''$ complectetur,
  et arcus super eam, qui est quantitas anomaliae,
  $1^\circ1'24''$\NM{8}: arcus igitur KC eclipticae,
  $31^\circ1'24''$\NM{9} circiter erit.
  Haec de anomalia sufficiunt, quam explicare volebamus.
\end{minipage}

\A{AB01-PDF0136-P03}
Dicit [auctor]: Hunc motum ita ad singulos gradus in tabulis
a puncto apogei posuimus; similiter reperiuntur aequatio simplex

\begin{multicols}{2}
\footnotesize

\A{AB01-PDF0136-N01}
\NN{1}{Secundae apud Platonem $55''$.}

\A{AB01-PDF0136-N02}
\NN{2}{Cod. « unius partis et trium ac triginta secundarum ».}

\A{AB01-PDF0136-N03}
\NN{3}{Cod. AD. Postea Plato 120.}

\A{AB01-PDF0136-N04}
\NN{4}{Cod. HDG.}

\A{AB01-PDF0136-N05}
\NN{5}{Cod. HS.}

\A{AB01-PDF0136-N06}
\NN{6}{Secundae apud Platonem $55''$.}

\A{AB01-PDF0136-N07}
\NN{7}{Plato $58^{p}15'34''$.
Cfr. adnot. \textit{e} huius capitis.}

\A{AB01-PDF0136-N08}
\NN{8}{Cod. $1^\circ4'24''$; Plato $1^\circ4'54''$.}

\A{AB01-PDF0136-N09}
\NN{9}{Cod. $1^\circ1'24''$; Plato $31^\circ4'54''$.}

\end{multicols}

% ================================================================
% PHYSICAL PDF137 — PRINTED PAGE 48
% ================================================================

\SP{0137}{48}{C. A. NALLINO}

\A{AB01-PDF0137-P01}
Lunae et aequatio media quinque planetarum, quae est
semidiametrus eorum epicycli, cum eius sinus sumatur et eo
modo dividatur.

\A{AB01-PDF0137-P02}
Si id calculo reperire vis, gradus epicycli quos ab apogeo
Sol, vel Luna, vel planeta percurrerit, scilicet anomaliam,
observa; si gradus minus quam $180^\circ$ sunt, iis utere,
sed si $180^\circ$ superant, eos de $360^\circ$ deme et
residuum adhibe. Graduum alterutra via repertorum, si minus
quam $90^\circ$ sunt, sinum atque cosinum multiplica in
semidiametrum epicycli sideris qui est totius aequationis
sinus; ambo producta per semidiametrum divide, et quod ex
divisione cosinus provenerit adde $60^{p}$ scilicet
semidiametro. Quadratum huius summae adde quadrato eius
quod e divisione sinus exierit; huius summae radicem inveni.
Deinde quod e divisione sinus provenerat, in semidiametrum
multiplica, et productum per inventam radicem divide.
— Si contra gradus adhibendi superant $90^\circ$, de iis
$90^\circ$ deme; sinum et cosinum residui in semidiametrum
epicycli multiplica et per semidiametrum divide; quod e
sinu provenerit de $60^{p}$ deme; quadratum residui adde
quadrato eius quod e cosinu exierit, et huius summae
radicem inveni. Postea quod e divisione cosinus provenerat
in semidiametrum multiplica, et productum per radicem
inventam divide. — Quod ex alterutra operatione tibi
provenerit, in arcum converte; arcus erit portio graduum
ad anomaliam qua usus es pertinens, scilicet aequatio
stellae \textit{(f)}.

\A{AB01-PDF0137-P03}
Semidiametrus epicycli Solis est $2^{p}4'45''$\NM{1};
semidiametrus epicycli Lunae $5^{p}15'$\NM{2};
Saturni $6^{p}29'50''$\NM{3}; Iovis $11^{p}30'0''$\NM{4};
Martis $39^{p}27'22''$\NM{5}; Veneris $43^{p}9'0''$\NM{6};
Mercurii $22^{p}30'30''$\NM{7}.
Hoc ex observationibus patuit et cum calculis congruit;
idque, volente Deo, est sinus aequationis mediae \textit{(g)}.

\A{AB01-PDF0137-H01}
\begin{center}
\textbf{CAPUT XXIX.}\par
De inaequalitate nychthemerōn et de aliorum in alia conversione.
\end{center}

\A{AB01-PDF0137-P04}
Dicit [auctor]: Plurimi homines et vulgus nychthemera
existimant aequalia esse, et eorum quodque 24 horas amplecti;
sed hoc haud verum est, nam nychthemerum medium est ortus
omnium 360 temporum aequatoris ab horizonte vel a circulo
meridiano, eo addito quod ex aequatoris temporibus ascendit
cum 59 minutis quae Sol, motu suo medio, in nychthemero
percurrit. Nychthemerum autem inaequale est ortus 360
temporum aequatoris, addito motu inaequali Solis in
nychthemero, qui necessarie minor aut maior est quam $59'$.

\A{AB01-PDF0137-P05}
Quoniam omni loco initium a circulo horizontis variat prout
in illo variant ascensiones signorum, initium autem ab hora
meridiana immobile est neque mutatur, ob aequalitatem
ascensionum signorum in circulo meridiano quolibet loco,
in calculo stellarum earumque locorum initium diei non
ponitur ab ortu vel occasu Solis, sed a tempore meridiei
aut mediae noctis\NM{8} \textit{(a)}.

\A{AB01-PDF0137-P06}
Item quoniam omnes motus stellarum in tabulis nonnisi ad
aequales dies supputantur, differentia inter nychthemera
inaequalia et nychthemera media negligitur. In motu Solis
et reliquarum stellarum differentia tanta non est ut
errorem sensibilem gignat; sed error mani-

% Sentence continues on PDF138 / printed page 49.
% The missing continuation is not supplied here.

\begin{multicols}{2}
\footnotesize

\A{AB01-PDF0137-N01}
\NN{1}{Cod. $5'$ pro $4'$.}

\A{AB01-PDF0137-N02}
\NN{2}{Gradus in cod. $6^\circ$.}

\A{AB01-PDF0137-N03}
\NN{3}{Secundae apud Plat. $7''$.}

\A{AB01-PDF0137-N04}
\NN{4}{Secundae apud Plat. $5''$.}

\A{AB01-PDF0137-N05}
\NN{5}{Cod. $19^\circ25'22''$; minuta apud Plat. $55'$.}

\A{AB01-PDF0137-N06}
\NN{6}{Plat. $44^\circ9'5''$.}

\A{AB01-PDF0137-N07}
\NN{7}{Gradus in cod. $29^\circ$.}

\A{AB01-PDF0137-N08}
\NN{8}{\textit{Almag.} III, 8 (ed. Halma, t. I, p. 208).}

\end{multicols}

\end{document}